% Chipathon-only supplement (3 frames). Gated by \ifchipathon in
% main.tex; built as standalone PDF + PNG export by main-chipathon.tex.

\begin{frame}{Track D digital: idea-to-chip on GF180MCU}
  \driving{LibreLane + cocotb GL sim}
  \footnotesize
  \begin{itemize}
    \item Natural-language spec to GDS in two phases:
          \cmd{idea\_loop} (sim-in-the-loop RTL) then
          \cmd{autoresearch} (greedy flow-knob optimisation).
    \item Every iteration is a real LibreLane flow plus
          post-synth and post-PnR-SDF cocotb GL sim.
          \warn{No \cmd{--stop-after} shortcut.}
    \item Same loop, two CLI backends (Claude Code, opencode).
          The proposer changes; the gate stays the same.
  \end{itemize}
  \vspace{0.6em}
  \begin{center}
  \fbox{\parbox{0.85\textwidth}{%
    \scriptsize\centering
    Live demo case:\\[0.2em]
    \cmd{python examples/11\_idea\_to\_chip\_demo\_gf180.py
         --case idea\_to\_optimize}\\
    \cmd{\hspace{2em} --backend \{cc\_cli|opencode\} --budget 3}
  }}
  \end{center}
  \note{%
    Chipathon supplement frame A. Mirror of the analog Track D
    slide. Voice-over: this is where the digital loop spends its
    Track D time. Two phases, same backend choice as the analog
    wizard.}
\end{frame}

\begin{frame}{Backend portability via the harness contract}
  \begin{center}
  \begin{tikzpicture}[node distance=0.6cm and 1.0cm,
                      every node/.style={font=\footnotesize}]
    \node[agent3, minimum width=4.0cm, line width=0.8mm] (loop)
      {Autoresearch loop\\\scriptsize \cmd{\_propose\_params}};
    \node[bxclaude,  below left =0.9cm and 0.5cm of loop,
          minimum width=3.4cm] (cc)
      {Claude Code CLI\\\scriptsize subscription};
    \node[bxopencode, below right=0.9cm and 0.5cm of loop,
          minimum width=3.4cm] (oc)
      {opencode CLI\\\scriptsize any provider};
    \draw[arrowclaude] (loop) -- (cc);
    \draw[arrowopen]   (loop) -- (oc);
  \end{tikzpicture}
  \end{center}
  \footnotesize
  \begin{itemize}
    \item Single async \cmd{run() $\rightarrow$ HarnessResult}
          contract behind both CLIs.
    \item Backend choice is per-user (subscription, provider
          preference); the loop body never branches.
    \item Same dispatch pattern as the analog wizard
          (\cmd{idea\_to\_rtl.py:\_build\_harness}).
  \end{itemize}
  \note{%
    Chipathon supplement frame B. Compact backends slide. Point
    out that the same architectural pattern holds across analog
    (sizing wizard) and digital (autoresearch) loops. The
    audience walks away knowing they can run Track D with whatever
    LLM subscription they already pay for.}
\end{frame}

\begin{frame}{Live A/B running today on GF180}
  \footnotesize
  \begin{center}
  \renewcommand{\arraystretch}{1.25}
  \begin{tabular}{@{}p{0.20\textwidth}p{0.37\textwidth}p{0.37\textwidth}@{}}
    \toprule
     & \textbf{\color{ccclaude}cc\_cli} & \textbf{\color{ccopencode}opencode} \\
    \midrule
    Backend  & Claude Code CLI (Opus 4.7) & opencode CLI \\
    Model    & subscription default & \cmd{openai/gpt-5.3-codex} \\
    Design   & Goertzel FP32, GF180MCU & Goertzel FP32, GF180MCU \\
    Budget   & 3 evals (\cmd{flow}) & 3 evals (\cmd{flow}) \\
    Outputs  & \cmd{phase\_optimize/results.tsv} & \cmd{phase\_optimize/results.tsv} \\
             & 3 evolution plots (PNG) & 3 evolution plots (PNG) \\
    \bottomrule
  \end{tabular}
  \end{center}
  \vspace{0.4em}
  \begin{center}
  \fbox{\parbox{0.85\textwidth}{%
    \footnotesize\centering
    Smoke-budget A/B. Numbers fit in a two-line table here when
    the runs close; we share the comparison and the workdirs at
    the next weekly.}}
  \end{center}
  \note{%
    Chipathon supplement frame C. Honest framing: today is a
    smoke-budget; the comparison is the experimental design, not
    the punchline. Real numbers come at the next weekly. Mention
    that the workdirs are reproducible artefacts, not just
    screenshots.}
\end{frame}
